% GENERATED from content/data/enrichment.json by mxbuild.enrichment; do not edit.
\mxunit{enrichment-invariants-notes}{Parité et invariants}
\mxsubtitle{Seconde · raisonnement \textperiodcentered{} Leçon et méthode \textperiodcentered{} 01}

\begin{prerequisites}
Entiers pairs et impairs ; preuve par contradiction.
\end{prerequisites}

\begin{keypoint}[Idées et méthode]
Un invariant est une quantité ou une propriété conservée par chaque opération autorisée. Si l’état visé n’a pas la même valeur de cet invariant que l’état initial, il est inaccessible.

Pour étudier la parité, on peut remplacer les entiers par leurs restes 0 ou 1 modulo 2. Deux nombres impairs ont une somme paire. Ajouter un nombre pair ne change pas la parité.

Méthode : décrire précisément un coup, chercher ce qu’il conserve, calculer l’invariant initial et celui de la cible. Une obstruction prouve une impossibilité ; l’absence d’obstruction ne prouve pas la possibilité.
\end{keypoint}

\begin{example}
On part de 0 et chaque coup ajoute ou retranche 6. Tous les nombres obtenus sont pairs. Atteindre 15 est donc impossible. Atteindre 14 est aussi impossible, mais la seule parité ne suffit pas à le prouver : il faut remarquer que tout résultat est multiple de 6.
\end{example}

\begin{activity}[À explorer]
Écrire quelques suites de coups, puis classer les nombres atteints. Trouver une propriété commune plus précise que la parité.
\end{activity}

